\section{Dependencies, scope, and further reading}
\label{a:dependencies}

\subsection{What was proved from what}

The ordinary comparison has a single dependency chain:
\[
\begin{tikzcd}[column sep=small,row sep=large]
\text{finite limits and regular images}\arrow[d]\\
\text{associative relation calculus}\arrow[d]\\
\text{Mal'cev difunctionality}\arrow[d]\\
\text{effective support quotients}\arrow[d]\\
\text{saturated descent and refinement cells}\arrow[d]\\
\text{Schreier refinement}\arrow[d]\\
\text{simple-factor rigidity}\arrow[d]\\
\text{Jordan--H\"older}\arrow[d]\\
\text{basis size and additive invariants}.
\end{tikzcd}
\]
Exactness enters at the effective quotients. Pointedness supplies the zero
classes used to concatenate cells into a series. The Mal'cev condition is
used where it is needed: to make every relation difunctional. Ordinary
compatible iterated quotients require regularity and exactness, but do not
require Mal'cev. The written proofs keep these dependencies separate.

The precise ordinary theorem is Theorem~\ref{j:jordan-holder}: in a
pointed Barr-exact Mal'cev category, any two given finite composition
series, retaining their regular quotient maps, have the same factors up
to permutation and isomorphism. No global chain condition is needed to
compare two finite series already supplied. This is a sufficient common
framework, not a claim that these are the logically weakest possible
axioms for every formulation of Jordan--H\"older.

The existence of a finite series is a separate issue. The basis flags
provide the required series for a vector space with a given finite basis.
Later, finite-dimensional module theory uses the now established dimension
function to prove existence of composition series. That later argument
does not feed back into the proof of dimension. One can define $K_0$ by
generators and exact-sequence relations before proving Jordan--H\"older.
What requires the factor comparison is the identification of $K_0$ with
the free abelian group on the simple isomorphism classes. Using that
freeness to prove factor uniqueness would reverse the dependency.

The source of associativity changes as the constructions develop. For the
original algebra it is an axiom equating two composites. For relation
composition, Baer addition, and Hall multiplication, it is a theorem:
both bracketings are obtained from the same complete data of witnesses,
tuples, or flags. For coherent quotients, the same universal-property
principle controls the higher compatibility of successive constructions.
These assertions specify the maps or data that do the work; a resemblance
between diagrams alone would not establish them.

\subsection{Three different comparison claims}

The scope of the page calculations is worth recording alongside the core
theorem. The categorical rectangle works for all objects and finite normal
series satisfying the exact Mal'cev hypotheses. The explicit nonabelian
page construction in Section~\ref{p:chapter} uses composition-series
input in groups with endomorphism operators; that simplicity is used to
prove the normality of differential images. The filtered module calculus
in Section~\ref{s:chapter} allows arbitrary finite filtrations and an
arbitrary filtration-preserving module map. Its convergence records
kernel and cokernel layers, and its staggered identity specialization
pairs simple factors.

Thus the nonabelian page calculation and the module spectral calculation
are compatible where both apply. Neither is asserted to provide a general
spectral sequence for arbitrary filtered groups. The categorical theorem
already supplies the unified Jordan--H\"older statement independently of
those more specialized calculations.

\subsection{The additional hypotheses behind reconstruction}

The universal additive record applies to composition-finite objects, with
the set-size qualification stated in Corollary~\ref{j:universal-monoid}.
For finite-length modules the category is also closed under submodules
and quotients, as explicitly proved. We did not assume that extra closure
for every exact Mal'cev category. Additive records ignore whether an
extension splits; the marked extension group retains its endpoint maps.

The probe theorem reconstructs finite-dimensional \emph{nilpotent}
$k[t]$-modules. Its second-difference formula is proved using canonical
subquotients, and the chain decomposition is constructed before the
reconstruction conclusion is drawn. The same tests fail for an invertible
operator, as the example shows. Hall multiplication uses finite fields
and finite-dimensional nilpotent modules so that every coefficient and
every basis expansion involved is finite. Its associativity counts the
same flags in two orders. Marked extension classes, isomorphism classes
of middle objects, and actual submodules are kept distinct.

\subsection{Where the higher foundations enter}

The higher part has its own declared inputs. Rigidification identifies
product-preserving theory models in spaces with a familiar simplicial
presentation. The Kan loop-group equivalence, loop--delooping recognition,
and Dold--Kan explain the group and module
instances. Effective groupoid quotients and universal colimits make
homotopy quotients compatible with base change. Straightening interprets
families over $BG$ as coherent actions. Stabilization and the derived
category supply mapping spectra, hearts, and fiber/cofiber sequences.
The tangent/cotangent theorems supply linear coefficients for algebras,
the representing object for derivations, and its transitivity sequence.

Those results are explicitly imported from the cited foundational
literature; this exposition does not reprove higher topos theory or the
construction of its model categories. Its deductions include the
double-coset homotopy pullback, the coefficient category over $BG$, the
marked extension classification and its automorphisms, and explicit
derived tensor and cotangent calculations. The ordinary internal-linear
equivalence $\operatorname{Ab}(\mathbf{Grp}/G)\simeq
\mathbb Z[G]\text{-}\mathbf{Mod}$ is proved directly.

The term ``higher'' does not imply an unrestricted higher
Jordan--H\"older theorem. Normality data and compatibility remain part of
the problem, and literal subobject chains from the trivial higher group
force discreteness by Proposition~\ref{h:monos-discrete}. What the higher
framework contributes here is a precise enlargement of the retained data
and several new reconstruction tools. Its ordinary discrete models
recover the common theorem when its exact Mal'cev hypotheses are checked.

\subsection{Audit and reproducibility}

The \href{https://tensorgami.com/notes/quotients-and-reconstruction-audit.html}{public audit report}
inventories all theorem-like statements, their proof or imported-foundation
status, the diagrams, references, and edition fingerprints. It distinguishes
the full review of the preceding edition from the subsequent check of these
small clarifications and the matching web, PDF, and source files.
The ordinary core and the higher chapters received
independent mathematical rereads; integration checks tracked changed
action conventions, quotient maps, normality, degrees, and dependencies.
Supplementary finite group and finite field calculations test examples
of the actual maps and counts. Such computations do not establish the
general theorems. Compilation and inspection of the rendered pages are
separate presentation checks, not proof certification. No proof assistant
formalization or historical priority claim is made.

The earlier edition is preserved separately. This edition reorganizes
the core around relations and exactness, so that the group and module
comparison theorems are genuine instances. It retains the explicit
extension and reconstruction calculations, and places the new higher
constructions after the common ordinary proof.

\begin{thebibliography}{99}
\addcontentsline{toc}{section}{References}
\bibitem{Knop}
Friedrich Knop, \emph{Tensor envelopes of regular categories},
Advances in Mathematics \textbf{214} (2007), 571--617.
\href{https://arxiv.org/pdf/math/0610552}{Author preprint};
Section 5 and Appendices A--B, especially Appendix B2.

\bibitem{HTT}
Jacob Lurie, \emph{Higher Topos Theory}, Princeton University Press, 2009.
\href{https://www.math.ias.edu/~lurie/papers/HTT.pdf}{Author PDF}.
See Corollary 5.5.9.3, Theorem 6.1.0.6, and Sections 2.2 and 3.2.

\bibitem{HA}
Jacob Lurie, \emph{Higher Algebra}, September 18, 2017.
\href{https://www.math.ias.edu/~lurie/papers/HA.pdf}{Author PDF}.
See Sections 1.3, 1.4.2--1.4.3, 5.2.6 and 7.3;
Theorem 7.1.2.13; Proposition 7.1.4.11;
Corollaries 7.3.3.6 and 7.3.4.14.

\bibitem{Badzioch}
Bernard Badzioch, \emph{Algebraic theories in homotopy theory},
Annals of Mathematics \textbf{155} (2002), 895--913.
\href{https://annals.math.princeton.edu/2002/155-3/p06}{Journal article};
\href{https://arxiv.org/abs/math/0110101}{Author preprint}.

\bibitem{GoerssJardine}
Paul G. Goerss and John F. Jardine, \emph{Simplicial Homotopy Theory},
Progress in Mathematics \textbf{174}, Birkh\"auser, 1999.
\href{https://www.sas.rochester.edu/mth/sites/doug-ravenel/otherpapers/Goerss-Jardine2.pdf}{Book PDF};
Chapter V, Proposition 6.3 and Corollary 6.4.

\bibitem{BeardsleyFox}
Jonathan Beardsley and Landon Fox, \emph{Higher Groups and Higher Normality},
\href{https://arxiv.org/html/2407.21210v3}{arXiv:2407.21210v3},
August 13, 2026. In particular Definition 4.1 and Section 5,
including the compatibility condition in Theorem 5.14.

\bibitem{Stacks}
The Stacks Project Authors, \emph{The Stacks Project},
\href{https://stacks.math.columbia.edu/}{online reference}.
The specific tags used are
\href{https://stacks.math.columbia.edu/tag/010I}{010I} (extensions),
\href{https://stacks.math.columbia.edu/tag/012K}{012K} (filtered complexes),
\href{https://stacks.math.columbia.edu/tag/019G}{019G} (Dold--Kan),
\href{https://stacks.math.columbia.edu/tag/06XP}{06XP} (Ext),
\href{https://stacks.math.columbia.edu/tag/0643}{0643} (projective resolutions),
\href{https://stacks.math.columbia.edu/tag/08QX}{08QX} (cotangent transitivity),
and \href{https://stacks.math.columbia.edu/tag/08SH}{08SH} (regular quotients).

\bibitem{Deligne}
Pierre Deligne, \emph{Th\'eorie de Hodge II},
Publications Math\'ematiques de l'IH\'ES \textbf{40} (1971), 5--57.
\href{https://www.numdam.org/item/10.1007/BF02684692.pdf}{Article PDF};
Section 1.2.1 discusses interchange of associated gradings.

\bibitem{Grandis}
Marco Grandis, \emph{Homotopy spectral sequences},
\href{https://arxiv.org/pdf/1007.0632}{arXiv:1007.0632}, 2010.
Sections 2--4 give a broader nonabelian exact-couple setting.
\end{thebibliography}
